\begin{definition}[Concrete zero BHist measure surface]
\label{def:measure-zero-bhist-surface}
The concrete zero $\mathsf{BHist}$ measure surface is the following
$\MeasureUp$ carrier instance over the BEDC history substrate. Its carrier row is
$$
\mathsf{ZeroMeasureCode}(h)\Longleftrightarrow \hsame(h,\emp).
$$
Its carrier classifier is
$$
h\sim_0 k \Longleftrightarrow \hsame(h,k).
$$
The displayed event inventory contains exactly the singleton empty-event row:
an event presentation is measurable precisely when its visible
$\mathsf{BHist}$ event code is classified with $\emp$. The measure endpoint is
the constant $\RealUp$ zero endpoint,
$$
\mu_0(E):=0_R,
$$
for every displayed measurable event $E$. The countable-sum endpoint for every
measurable sequence is also the constant endpoint $0_R$.

The only ambient rows used by this surface are $\mathsf{BHist}$,
$\hsame$, the classifier and zero row of $\NameCert_{\RealUp}$, and the
additive $\AbGroupUp$ row of \autoref{def:abgroup-carrier} for zero, addition
congruence, and the additive zero laws.
\end{definition}

\begin{theorem}[Zero BHist measure carrier-classifier stability]
\label{thm:measure-zero-bhist-carrier-classifier-stability}
The concrete zero surface of
\autoref{def:measure-zero-bhist-surface} satisfies the carrier transport,
event-row transport, and endpoint transport rows required by the $\MeasureUp$
carrier classifier. If $\mathsf{ZeroMeasureCode}(h)$ and $h\sim_0 k$, then
$\mathsf{ZeroMeasureCode}(k)$. If $E$ is a displayed measurable event over
$h$ and $E'$ is transported along the same $\hsame$ comparison, then $E'$ is
again the singleton empty-event row over $k$. The endpoint comparison is
$$
\mu_0(E)\sim_R\mu_0(E').
$$
\end{theorem}
\begin{proof}
Unpack the carrier row. The assumption
$\mathsf{ZeroMeasureCode}(h)$ identifies $h$ with $\emp$, while $h\sim_0 k$
is exactly an $\hsame(h,k)$ comparison. Symmetry and transitivity of
$\hsame$ identify $k$ with $\emp$, so the transported carrier row is again
$\mathsf{ZeroMeasureCode}(k)$.

For event-row transport, the only displayed measurable row is the empty-event
row. Transporting an empty event presentation along an $\hsame$ comparison
therefore has only one possible visible target row: the same empty-event row
read in the $k$-fiber. No quotient of events or external set equality is used.

For endpoint transport, both endpoint projections reduce by definition to
$0_R$. The reflexive classifier row of $\NameCert_{\RealUp}$ classifies
$0_R$ with itself, giving $\mu_0(E)\sim_R\mu_0(E')$.
\end{proof}
\leanchecked{BEDC.Derived.MeasureUp.MeasureZeroBHist\_carrier\_classifier\_stability}

\begin{theorem}[Zero BHist measure sigma-additivity]
\label{thm:measure-zero-bhist-sigma-additivity}
For the concrete zero surface of
\autoref{def:measure-zero-bhist-surface}, every displayed pairwise-disjoint
measurable sequence $E_\bullet$ has countable union classified with the same
singleton empty-event row, and the measure comparison is
$$
\mu_0\left(\bigcup_{n:\mathbb{N}}E_n\right)\sim_R
\sum_{n=0}^{\infty}\mu_0(E_n).
$$
Here the countable-sum endpoint on the right is the constant endpoint $0_R$.
The head-tail readback for this endpoint is classified by the additive
$\AbGroupUp$ zero laws.
\end{theorem}
\begin{proof}
Every displayed measurable term in the zero surface is the empty-event row.
Thus each $E_n$ is classified with $\emp$, and the visible countable union of
the displayed sequence is again classified with $\emp$. The pairwise
disjointness premise adds no extra branch, because the only event row on either
side of an intersection is empty.

The measure endpoint of the union is $0_R$ by definition. Each summand
$\mu_0(E_n)$ is also $0_R$, and the countable-sum endpoint selected by the
surface is the constant $0_R$. Reflexivity of the $\RealUp$ classifier supplied
by $\NameCert_{\RealUp}$ therefore gives the displayed sigma-additivity
comparison.

It remains to read the head-tail row. The head endpoint is $0_R$, and the tail
endpoint is again $0_R$. The additive $\AbGroupUp$ row classifies
$0_R+0_R$ with $0_R$, and classifier symmetry gives the required comparison
between the constant sum endpoint and its head-tail expression.
\end{proof}

\begin{theorem}[Zero BHist MeasureUp NameCert package]
\label{thm:measure-zero-bhist-namecert-package}
The concrete zero surface of
\autoref{def:measure-zero-bhist-surface}, together with
\autoref{thm:measure-zero-bhist-carrier-classifier-stability} and
\autoref{thm:measure-zero-bhist-sigma-additivity}, forms a concrete
$\MeasureUp$ NameCert obligation package by the packaging theorem
\autoref{thm:measure-namecert-obligation-package}.
\end{theorem}
\begin{proof}
Project the concrete surface rows. The carrier row is
$\mathsf{ZeroMeasureCode}$, the classifier row is $\sim_0$, the measurable-event
inventory is the singleton empty-event row, the measure endpoint is the constant
$0_R$ map, and the countable-sum endpoint is the constant $0_R$ row.

The carrier-classifier theorem supplies the transport fields required by the
package: carried histories move along $\hsame$, measurable rows remain the
empty-event row, and endpoints compare by the reflexive $\RealUp$ zero
classifier. The sigma-additivity theorem supplies the countable-union row and
the head-tail readback row, with additive compatibility read from the
$\AbGroupUp$ zero laws.

These are exactly the rows consumed by
\autoref{thm:measure-namecert-obligation-package}. Repacking them through that
theorem gives a concrete $\MeasureUp$ obligation package whose visible carrier
is the zero $\mathsf{BHist}$ surface.
\end{proof}

\begin{theorem}[Zero BHist relative-difference zero row]
\label{thm:measure-zero-bhist-relative-difference-zero-row}
For the concrete zero $\mathsf{BHist}$ measure surface of
\autoref{def:measure-zero-bhist-surface}, let $h$ be a carried history with
$\mathsf{ZeroMeasureCode}(h)$. Let $A$ and $B$ be displayed measurable events
in the $h$-fiber, and suppose the public event readback displays
$B\subseteq A$. Then the zero surface realizes the
complement/relative-difference row for this inclusion by the singleton
empty-event row:
$$
A\setminus B\quad\text{is displayed by the event code }\emp.
$$
Consequently $A\setminus B$ is accepted as a displayed measurable event of the
zero surface, the row supplies the disjointness witness for $B$ and
$A\setminus B$ and the binary-union readback classifying
$B\cup(A\setminus B)$ with $A$, and the visible endpoint is classified as
zero:
$$
\mu_0(A\setminus B)\sim_R0_R.
$$
\end{theorem}
\begin{proof}
Unpack \autoref{def:measure-zero-bhist-surface}. The measurable-event inventory
of the zero surface has exactly one displayed row: an event presentation is
measurable precisely when its visible $\mathsf{BHist}$ event code is
classified with $\emp$. Since $A$ and $B$ are displayed measurable events in
the same zero fiber, both of their visible event codes are therefore classified
with $\emp$.

Assign the relative-difference output for the displayed inclusion
$B\subseteq A$ to the same singleton empty-event row. This is accepted by the
measurable-event inventory just read. The disjointness witness is the empty-row
intersection of $B$ with that singleton output: the $B$ input is already
classified with $\emp$, and the output is classified with $\emp$.

The binary-union readback is obtained by reading the union of two displayed
empty rows as the empty row and then composing, by symmetry and transitivity of
the zero event classifier, with the displayed comparison classifying $A$ with
that same row. This realizes the concrete complement/relative-difference field
described in \autoref{def:measure-complement-difference-row} for the zero
surface.

The endpoint projection of \autoref{def:measure-zero-bhist-surface} is
constant: for every displayed measurable event $E$, $\mu_0(E):=0_R$. Applying
this to the accepted output $A\setminus B$ gives the visible endpoint $0_R$.
Reflexivity of the $\RealUp$ classifier supplied by $\NameCert_{\RealUp}$
classifies $0_R$ with itself, hence
$$
\mu_0(A\setminus B)\sim_R0_R.
$$
\end{proof}
